% Synthetic Beamer-like slides (454×255 pt) for tests/test_extract_slides.py, compiled with
% XeLaTeX through h0lon.render.latex.compile_latex; the PDF next to it is committed.
% Slide 1: title + bullets ($\bullet$ from CMSY, like Beamer); slide 2: formulas;
% slide 3: a two-line title and text; «N/3» in the bottom right corner of every slide.
\documentclass{article}
\usepackage[paperwidth=454pt,paperheight=255pt,left=14pt,right=14pt,top=8pt,bottom=4pt]{geometry}
\usepackage{fontspec}
\IfFontExistsTF{Arial}{\setmainfont{Arial}}{\IfFontExistsTF{DejaVu Sans}{\setmainfont{DejaVu Sans}}{}}
\usepackage{amsmath}
\pagestyle{empty}
\setlength{\parindent}{0pt}
\newcommand{\slidetitle}[1]{{\fontsize{14.3}{17}\selectfont\bfseries #1\par}\vspace{14pt}}
\newcommand{\slidenum}[1]{\vfill{\hfill\fontsize{10}{12}\selectfont\bfseries #1/3}}
\begin{document}

\slidetitle{Метрические методы}
\begin{itemize}
  \item[$\bullet$] Объекты одного класса близки друг к другу.
  \item[$\bullet$] Похожесть задаётся расстоянием между объектами.
  \item[$\bullet$] Классификация по ближайшим соседям.
\end{itemize}
\slidenum{1}
\newpage

\slidetitle{Расстояние Минковского}
\[
  \rho(a, b) = \Bigl(\sum_{i} |a_i - b_i|^p\Bigr)^{1/p}
\]
\begin{itemize}
  \item[$\bullet$] При $p = 1$ это манхэттенское расстояние.
  \item[$\bullet$] При $p = 2$ это евклидово расстояние $\sqrt{\sum_i (a_i - b_i)^2}$.
\end{itemize}
\slidenum{2}
\newpage

\slidetitle{Ядерное сглаживание\\и непараметрическая регрессия}
Оценка плотности строится по окну вокруг точки.
Ширина окна определяет гладкость оценки.
\slidenum{3}

\end{document}
